% Example document using hydra-doc.sty
% Compile from docs/templates/:  pdflatex example.tex && pdflatex example.tex
\documentclass[11pt,a4paper]{article}
\usepackage{hydra-doc}

%% Logo is in the same directory when compiling from templates/
\setlogopath{newpaltz-logo}

%% Header text
\fancyhead[L]{\small Example Document}
\fancyhead[R]{\small SUNY New Paltz}

\hydradate{February 2026}

\begin{document}

\hydratitlepage{Example Document}{Hydra Doc Template Demo}{Infrastructure Team}

\tableofcontents
\newpage

%% ============================================================
\section{Getting Started}
%% ============================================================

To use the SUNY New Paltz CS document template, add this to your preamble:

\begin{lstlisting}[style=bash]
\usepackage{hydra-doc}        % when compiling from docs/templates/
\usepackage{templates/hydra-doc}  % when compiling from docs/
\end{lstlisting}

The package loads all required packages, defines official brand colors,
listing styles, callout box environments, and sets up headers/footers.

%% ============================================================
\section{Brand Colors}
%% ============================================================

The template uses the official SUNY New Paltz brand colors:

\begin{table}[H]
\centering
\begin{tabular}{@{}lll@{}}
\toprule
\textbf{Color} & \textbf{Name} & \textbf{Hex} \\
\midrule
\textcolor{np-blue}{\rule{1em}{1em}} & \texttt{np-blue} (Pantone 281) & \texttt{\#003E7E} \\
\textcolor{np-orange}{\rule{1em}{1em}} & \texttt{np-orange} (Pantone 165) & \texttt{\#F58426} \\
\textcolor{np-gold}{\rule{1em}{1em}} & \texttt{np-gold} (Pantone 7409) & \texttt{\#FDB924} \\
\textcolor{np-lightblue}{\rule{1em}{1em}} & \texttt{np-lightblue} (Pantone 639) & \texttt{\#00A5D9} \\
\textcolor{np-coral}{\rule{1em}{1em}} & \texttt{np-coral} (Pantone 7417) & \texttt{\#F26649} \\
\textcolor{np-lime}{\rule{1em}{1em}} & \texttt{np-lime} (Pantone 583) & \texttt{\#B0BC22} \\
\textcolor{np-slate}{\rule{1em}{1em}} & \texttt{np-slate} (Pantone 5425) & \texttt{\#80A1B6} \\
\textcolor{np-brown}{\rule{1em}{1em}} & \texttt{np-brown} (Pantone 174) & \texttt{\#A84D10} \\
\bottomrule
\end{tabular}
\caption{Official SUNY New Paltz brand colors}
\end{table}

%% ============================================================
\section{Callout Boxes}
%% ============================================================

Four callout box environments are available:

\begin{infobox}
\textbf{Info:} Use \texttt{infobox} for general information and tips.
\end{infobox}

\begin{warningbox}
\textbf{Warning:} Use \texttt{warningbox} for important warnings.
\end{warningbox}

\begin{successbox}
\textbf{Success:} Use \texttt{successbox} for confirmations and good outcomes.
\end{successbox}

\begin{dangerbox}
\textbf{Danger:} Use \texttt{dangerbox} for critical alerts.
\end{dangerbox}

Inline callout commands:

\notebox{This is an inline note created with \cmdref{\textbackslash notebox\{...\}}.}

\warnbox{This is an inline warning created with \cmdref{\textbackslash warnbox\{...\}}.}

%% ============================================================
\section{Code Listings}
%% ============================================================

Three listing styles are predefined: \texttt{bash}, \texttt{yaml}, and \texttt{js}.

\subsection{Bash}

\begin{lstlisting}[style=bash]
# Check cluster health
kubectl get nodes -o wide
kubectl get pods -A | grep -v Running
\end{lstlisting}

\subsection{YAML}

\begin{lstlisting}[style=yaml]
apiVersion: v1
kind: Service
metadata:
  name: hydra-auth
  namespace: hydra-system
spec:
  ports:
    - port: 6969
      targetPort: 6969
\end{lstlisting}

\subsection{JavaScript}

\begin{lstlisting}[style=js]
const express = require('express');
const app = express();

app.get('/health', (req, res) => {
    res.json({ status: 'ok' });
});
\end{lstlisting}

%% ============================================================
\section{Tables}
%% ============================================================

Use \texttt{booktabs} for clean tables:

\begin{table}[H]
\centering
\begin{tabular}{@{}llr@{}}
\toprule
\textbf{Node} & \textbf{Role} & \textbf{RAM} \\
\midrule
Hydra    & Control plane   & 256 GB \\
Chimera  & GPU inference   & 64 GB  \\
Cerberus & GPU training    & 64 GB  \\
\bottomrule
\end{tabular}
\caption{Cluster nodes}
\end{table}

%% ============================================================
\section{References and Links}
%% ============================================================

Use \cmdref{\textbackslash cmdref\{...\}} for inline monospace references
like \cmdref{kubectl get pods} or file paths like \cmdref{/etc/rancher/rke2/config.yaml}.

Hyperlinks work as expected: \url{https://docs.rke2.io/}

\end{document}
